\documentclass[11pt]{article}
\usepackage[margin=1in]{geometry}
\usepackage{amsmath,amssymb}
\usepackage[T1]{fontenc}
\title{An Analytical Balance Step for One Selected Voronoi Cell}
\author{Working research note: supplement to the recovered protocol}
\date{29 September 2026}
\begin{document}
\maketitle

\paragraph{Status.}
This supplement records the agreed ordinary balance step. It supersedes
the preceding combined-badness acceptance rule for balance moves in the
recovered research note. It does not replace that note's historical
experiments or its departure, migration, and insertion protocol.
No implementation or simulation of this revised balance rule is claimed.
Assume cell $i$ has already been selected; selection and scheduling remain
deferred.

\section{High-level rule}
\begin{enumerate}
\item Compute current cell densities and use them to predict weights as
center $i$ moves. All other centers remain fixed.
\item Measure imbalance by summing the mean absolute neighbor-weight gaps
for $i$ and its original neighbors.
\item Compute analytically the direction of fastest first-order decrease
of this neighborhood objective.
\item Compute a conservative safe distance analytically, certifying decrease
throughout the segment, preserving adjacency and valid geometry, keeping
the center in the domain, and avoiding collisions.
\item Move once by that distance and end the step. If no strictly decreasing
feasible direction exists, stay put. Do not execute an uncertified segment.
\end{enumerate}

\section{Computed densities and neighborhood objective}
Let $P$ be the fixed bounded convex polygonal domain. Let $p_i$ be the
initial position of the selected center and $q$ its candidate position.
Write $W_j^0$ and $A_j^0>0$ for each cell's current weight and area.
Compute
\[
 \rho_j=\frac{W_j^0}{A_j^0},\qquad
 \widehat W_j(q)=\rho_j A_j(q).
\]
The computed densities are parameters of this step's weight model.
Candidate weights are modeled weights, not recounted pickups. For an
ordinary count-based starting state, $W_j^0$ is the current pickup count;
weight accounting following departure or insertion remains a separate
implementation question in the preceding protocol.

Let $N_j^0$ contain the initial neighbors sharing a positive-length edge,
and fix $\mathcal S_i=\{i\}\cup N_i^0$. On a segment preserving adjacency,
define
\[
 \widehat G_j(q)=\frac{1}{|N_j^0|}\sum_{k\in N_j^0}
 |\widehat W_j(q)-\widehat W_k(q)|,
 \qquad E_i(q)=\sum_{j\in\mathcal S_i}\widehat G_j(q).
\]
An empty neighbor set is an invalid scoring state. This is a sum of
local mean gaps, not a median and not a combined balance--circularity
score. It protects the neighborhood aggregate, not every individual
score or the global median.

\section{Analytical direction}
For each initial shared edge of $i$ and $j$, let $\ell_{ij}$ denote its
length and $m_{ij}$ its midpoint. Define
\[
 b_{ij}=\frac{\ell_{ij}}{\|p_j-p_i\|}(m_{ij}-p_i),\qquad
 a_i=\sum_{j\in N_i^0}b_{ij},\qquad
 a_j=-b_{ij}\quad(j\in N_i^0).
\]
Set $a_j=0$ for other cells. In a nondegenerate configuration with fixed
domain boundary, $a_j=\nabla_{p_i}A_j$. Thus a small displacement $v$
has area differential $dA_j=a_j\cdot v$. Put $z_j=\rho_j a_j$.

If every compared weight difference is nonzero, the objective gradient is
\[
 g=\sum_{j\in\mathcal S_i}\frac{1}{|N_j^0|}
 \sum_{k\in N_j^0}\operatorname{sgn}(W_j^0-W_k^0)(z_j-z_k).
\]
For an interior center with $g\ne0$, the steepest direction per unit
distance is $u=-g/\|g\|$. If $g=0$, no strictly decreasing first-order
direction exists; this does not establish optimality against finite moves.

At equal weights, use the exact directional derivative
\[
 D E_i(u)=\sum_{j\in\mathcal S_i}\frac{1}{|N_j^0|}
 \sum_{k\in N_j^0}
 \begin{cases}
 \operatorname{sgn}(W_j^0-W_k^0)(z_j-z_k)\cdot u,
       &W_j^0\ne W_k^0,\\
 |(z_j-z_k)\cdot u|,&W_j^0=W_k^0.
 \end{cases}
\]
Choose a feasible unit direction minimizing this expression and require
a negative minimum. In two dimensions, the sign boundaries partition
the direction circle into finitely many arcs; on each arc the expression
is linear in $u$, allowing analytical candidate evaluation. Boundary
centers require directions admitting a positive segment inside $P$.
Tie-breaking between equally minimizing directions requires a
deterministic implementation convention.

\section{One analytically certified safe segment}
Set $f(t)=E_i(p_i+tu)$ and $a=-f'(0+)>0$. Obtain $R>0$ such that the
entire interval preserves valid geometry, the domain constraint,
noncollision, the current geometric structure, and the relevant
one-sided weight-order branches. Choose $R$ strictly before any positive
event that would invalidate these conditions. Obtain a bound $M\ge0$ with
\[
 |f''(t)|\le M\qquad(0<t\le R).
\]
Use
\[
 h=\begin{cases}\min\{R,a/(2M)\},&M>0,\\R,&M=0,\end{cases}
 \qquad p_i^{\mathrm{new}}=p_i+hu.
\]
For $0<t\le h$, $f'(t)\le-a+Mt\le-a/2<0$. Consequently,
\[
 E_i(p_i^{\mathrm{new}})\le E_i(p_i)-ah/2.
\]
This certifies decrease throughout the segment, rather than checking only
its endpoint. It is conservative and need not be the longest safe move.
The step ends after this single segment: there are no repeated direction
or distance recalculations within the selected cell's turn.

\subsection{Obtaining analytical bounds}
For a fixed polygonal domain and nondegenerate combinatorial structure,
Voronoi vertices are intersections of lines whose coefficients are
polynomial in $t$. Their coordinates, cell areas, and the objective on
a fixed absolute-value branch are rational functions of $t$.
Geometric event predicates and denominator conditions can therefore be
expressed as polynomial sign conditions.

A conservative sign-preserving radius can be computed without root
search. For $Q(t)=\sum_{k=0}^d c_k t^k$ with $c_0\ne0$, choose a
reference distance $T>0$ and set $B=\sum_{k=1}^d|c_k|T^k$. Use
\[
 r_Q=\begin{cases}
 T\min\{1,|c_0|/(2B)\},&B>0,\\T,&B=0.
 \end{cases}
\]
For $0\le t\le r_Q$, the nonconstant terms have total magnitude at most
$|c_0|/2$. A minimum of such radii over sufficient event and denominator
conditions gives a conservative $R$. When a polynomial vanishes at zero,
factor out its initial power of $t$ to determine the right-hand sign;
identically zero predicates require separate handling. Bounding the
rational second derivative using numerator bounds and denominators
bounded away from zero gives $M$. No trial movements or numerical
line search are required by this construction.

\subsection{Scope of certification}
Equal-weight ties require the right-hand branches described above.
Initial geometric degeneracies require an explicit one-sided convention
before a positive segment can be certified. Enumerating sufficient event
predicates, deterministic conventions, and rigorous numerical bounds
remain implementation work. Ordinary floating-point estimates without
error bounds are not mathematical certificates; exact or outward-rounded
calculations are needed for a numerical guarantee.

Two-hop information suffices to evaluate the neighborhood objective,
given the relevant cell geometry and weights. This alone does not prove
that all geometric event certificates can be obtained from only two-hop
information. Additional local information exchange may be required.

\section{Separate circularity and deferred decisions}
The agreed intended mixture among ordinary moves is 90\% balance and
10\% circularity. Circularity steps should improve
$C_i=4\pi A_i/L_i^2$ without increasing $E_i$ throughout the move.
Their precise direction and certified-distance rule remains to be
specified. There is no combined badness or balance--circularity exchange
coefficient in the revised ordinary movement rule.

Cell selection, including whether to sample in proportion to current
$G_i$, replacement, and score-refresh timing, remains deferred. Earlier
badness-weighted selection is historical, not an adopted scheduler for
this revision. The computed per-cell densities can differ, so total
modeled weight need not be conserved as areas change. Descent of $E_i$
does not guarantee improvement of refreshed pickup-count imbalance or
convergence. The departure--migration--insertion protocol and its
remaining weight-accounting questions are unchanged.
\end{document}
